% Options for packages loaded elsewhere
\PassOptionsToPackage{unicode}{hyperref}
\PassOptionsToPackage{hyphens}{url}
\PassOptionsToPackage{dvipsnames,svgnames,x11names}{xcolor}
%
\documentclass[
  11pt,
]{article}
\usepackage{amsmath,amssymb}
\usepackage{iftex}
\ifPDFTeX
  \usepackage[T1]{fontenc}
  \usepackage[utf8]{inputenc}
  \usepackage{textcomp} % provide euro and other symbols
\else % if luatex or xetex
  \usepackage{unicode-math} % this also loads fontspec
  \defaultfontfeatures{Scale=MatchLowercase}
  \defaultfontfeatures[\rmfamily]{Ligatures=TeX,Scale=1}
\fi
\usepackage{lmodern}
\ifPDFTeX\else
  % xetex/luatex font selection
\fi
% Use upquote if available, for straight quotes in verbatim environments
\IfFileExists{upquote.sty}{\usepackage{upquote}}{}
\IfFileExists{microtype.sty}{% use microtype if available
  \usepackage[]{microtype}
  \UseMicrotypeSet[protrusion]{basicmath} % disable protrusion for tt fonts
}{}
\makeatletter
\@ifundefined{KOMAClassName}{% if non-KOMA class
  \IfFileExists{parskip.sty}{%
    \usepackage{parskip}
  }{% else
    \setlength{\parindent}{0pt}
    \setlength{\parskip}{6pt plus 2pt minus 1pt}}
}{% if KOMA class
  \KOMAoptions{parskip=half}}
\makeatother
\usepackage{xcolor}
\usepackage[margin=1in]{geometry}
\usepackage{color}
\usepackage{fancyvrb}
\newcommand{\VerbBar}{|}
\newcommand{\VERB}{\Verb[commandchars=\\\{\}]}
\DefineVerbatimEnvironment{Highlighting}{Verbatim}{commandchars=\\\{\}}
% Add ',fontsize=\small' for more characters per line
\newenvironment{Shaded}{}{}
\newcommand{\AlertTok}[1]{\textcolor[rgb]{1.00,0.00,0.00}{\textbf{#1}}}
\newcommand{\AnnotationTok}[1]{\textcolor[rgb]{0.38,0.63,0.69}{\textbf{\textit{#1}}}}
\newcommand{\AttributeTok}[1]{\textcolor[rgb]{0.49,0.56,0.16}{#1}}
\newcommand{\BaseNTok}[1]{\textcolor[rgb]{0.25,0.63,0.44}{#1}}
\newcommand{\BuiltInTok}[1]{\textcolor[rgb]{0.00,0.50,0.00}{#1}}
\newcommand{\CharTok}[1]{\textcolor[rgb]{0.25,0.44,0.63}{#1}}
\newcommand{\CommentTok}[1]{\textcolor[rgb]{0.38,0.63,0.69}{\textit{#1}}}
\newcommand{\CommentVarTok}[1]{\textcolor[rgb]{0.38,0.63,0.69}{\textbf{\textit{#1}}}}
\newcommand{\ConstantTok}[1]{\textcolor[rgb]{0.53,0.00,0.00}{#1}}
\newcommand{\ControlFlowTok}[1]{\textcolor[rgb]{0.00,0.44,0.13}{\textbf{#1}}}
\newcommand{\DataTypeTok}[1]{\textcolor[rgb]{0.56,0.13,0.00}{#1}}
\newcommand{\DecValTok}[1]{\textcolor[rgb]{0.25,0.63,0.44}{#1}}
\newcommand{\DocumentationTok}[1]{\textcolor[rgb]{0.73,0.13,0.13}{\textit{#1}}}
\newcommand{\ErrorTok}[1]{\textcolor[rgb]{1.00,0.00,0.00}{\textbf{#1}}}
\newcommand{\ExtensionTok}[1]{#1}
\newcommand{\FloatTok}[1]{\textcolor[rgb]{0.25,0.63,0.44}{#1}}
\newcommand{\FunctionTok}[1]{\textcolor[rgb]{0.02,0.16,0.49}{#1}}
\newcommand{\ImportTok}[1]{\textcolor[rgb]{0.00,0.50,0.00}{\textbf{#1}}}
\newcommand{\InformationTok}[1]{\textcolor[rgb]{0.38,0.63,0.69}{\textbf{\textit{#1}}}}
\newcommand{\KeywordTok}[1]{\textcolor[rgb]{0.00,0.44,0.13}{\textbf{#1}}}
\newcommand{\NormalTok}[1]{#1}
\newcommand{\OperatorTok}[1]{\textcolor[rgb]{0.40,0.40,0.40}{#1}}
\newcommand{\OtherTok}[1]{\textcolor[rgb]{0.00,0.44,0.13}{#1}}
\newcommand{\PreprocessorTok}[1]{\textcolor[rgb]{0.74,0.48,0.00}{#1}}
\newcommand{\RegionMarkerTok}[1]{#1}
\newcommand{\SpecialCharTok}[1]{\textcolor[rgb]{0.25,0.44,0.63}{#1}}
\newcommand{\SpecialStringTok}[1]{\textcolor[rgb]{0.73,0.40,0.53}{#1}}
\newcommand{\StringTok}[1]{\textcolor[rgb]{0.25,0.44,0.63}{#1}}
\newcommand{\VariableTok}[1]{\textcolor[rgb]{0.10,0.09,0.49}{#1}}
\newcommand{\VerbatimStringTok}[1]{\textcolor[rgb]{0.25,0.44,0.63}{#1}}
\newcommand{\WarningTok}[1]{\textcolor[rgb]{0.38,0.63,0.69}{\textbf{\textit{#1}}}}
\setlength{\emergencystretch}{3em} % prevent overfull lines
\providecommand{\tightlist}{%
  \setlength{\itemsep}{0pt}\setlength{\parskip}{0pt}}
\setcounter{secnumdepth}{5}
\usepackage{amsmath,amssymb,amsthm}
\ifLuaTeX
  \usepackage{selnolig}  % disable illegal ligatures
\fi
\usepackage{bookmark}
\IfFileExists{xurl.sty}{\usepackage{xurl}}{} % add URL line breaks if available
\urlstyle{same}
\hypersetup{
  pdftitle={Intrinsic Reeb comparison and signed cutoff estimates},
  pdfauthor={Source-only research draft},
  colorlinks=true,
  linkcolor={Maroon},
  filecolor={Maroon},
  citecolor={Blue},
  urlcolor={Blue},
  pdfcreator={LaTeX via pandoc}}

\title{Intrinsic Reeb comparison and signed cutoff estimates}
\usepackage{etoolbox}
\makeatletter
\providecommand{\subtitle}[1]{% add subtitle to \maketitle
  \apptocmd{\@title}{\par {\large #1 \par}}{}{}
}
\makeatother
\subtitle{A partial argument for unchanged P13.1}
\author{Source-only research draft}
\date{4 October 2026}

\begin{document}
\maketitle

\subsection{Scope and status}\label{scope-and-status}

This document proves an intrinsic comparison of the actual classes
\(\mathcal H(\xi)\), \(\operatorname{CPSH}(\xi)\) and
\(\operatorname{PSH}(\xi)\) for arbitrary, possibly irrational and
nonproportional, Reeb fields on every cone in the problem. It also
proves signed finite-type norm representation and quantitative signed
cutoff estimates. The latter are algebraic estimates, not estimates for
the genuine affine-cone energy. The required preservation of
\(\mathcal E^1\), control of moving mixed Monge--Ampère measures, and
uniform differentiation of the genuine extended energy are not proved.
Consequently this is not a solution of P13.1.

The eight supplied caches are the only mathematical sources used. Their
identifiers in this document are A (Wu, 2406.13684v1), B (Wu,
2106.13395v2), C (Boucksom--Jonsson, 1801.08229v4), D
(Boucksom--Jonsson, 2307.01697v1), E (Lyuhui Wu, 2609.27610v1), F
(Boucksom--Jonsson, 1805.11160v1), G (Chambert-Loir--Ducros,
1204.6277v2), and H (Boucksom--Hisamoto--Jonsson, 1504.06568v3). All
line numbers below are one-based inclusive ranges in the supplied UTF-8
extract split by the literal LF character. Source results used as
premises are identified explicitly. No earlier research note is a
premise.

\subsection{1. Fixed carrier, action, and
normalization}\label{fixed-carrier-action-and-normalization}

Let \(R\) be the coordinate ring of the normal cone, of dimension \(n\),
with the good torus action of the problem, \(R_0=k\), and homogeneous
maximal ideal \(\mathfrak m\). Choose nonzero homogeneous algebra
generators \(h_1,\ldots,h_q\in\mathfrak m\) with weights \(\beta_i\).
Write \(d_\xi(\alpha)=\langle\xi,\alpha\rangle\). Every supported
nonzero weight is a nonnegative integral combination of the \(\beta_i\):
take a homogeneous polynomial expression of a nonzero element of
\(R_\alpha\) and retain a nonzero monomial. In particular the generators
control all degree ratios, not just the weights in a particular FS
presentation.

We use the semivaluation convention \(\log|f|(v)=-v(f)\). Torification
is \(q(v)(\sum f_\alpha)=\min_\alpha v(f_\alpha)\). On its image let \[
 (v+t\eta)(f)=\min_\alpha\{v(f_\alpha)+t d_\eta(\alpha)\}.
 \tag{1}
\] These are the analytic action and torification recorded in A,
L218--246; hence they are semivaluations, not merely formal functions.
The action is additive on torified points. It is jointly continuous,
since evaluation on each polynomial is a minimum of finitely many
continuous homogeneous evaluations. This also proves continuity of
\(q\).

Set \(Y=q(X^{\rm an}\setminus\{o\})\) and \[
 K=X_0^T=\{v\in Y: v(\mathfrak m)=0\}.
\] This is one fixed carrier for all Reeb fields. On \(K\) all regular
functions have nonnegative value: a polynomial in the \(h_i\) has
nonnegative value when all \(v(h_i)\geq0\). Moreover
\(v(\mathfrak m)=\min_i v(h_i)\). The closed unit polydisc in the
coordinates \(h_i\) is compact; imposing the equations of \(R\) and
\(\min_i v(h_i)=0\) gives a closed compact subset. Imposing \(q(v)=v\)
is a further closed condition. Thus \(K\) is compact.

The canonical reference function and the orbit section are \[
 \rho_\xi(v)=\max_i\frac{-v(h_i)}{d_\xi(\beta_i)},\qquad
 s_\xi(v)=v+\rho_\xi(v)\xi\in K.                 \tag{2}
\] At least one \(h_i\) has finite value for \(v\ne o\), so the maximum
is finite. Each generator has nonnegative value at \(s_\xi(v)\), and at
least one has value zero. Also \(\rho_\xi(v+t\xi)=\rho_\xi(v)-t\).
Consequently a torified \(\xi\)-equivariant function \(u\) obeys \[
 u(v)=u(s_\xi(v))+\rho_\xi(v).                  \tag{3}
\] In particular its restriction to \(K\) determines it on \(Y\), and
then on \(X^{\rm an}\) by torification and the value \(-\infty\) at
\(o\). Uniform differences on \(K\) are uniform differences on all of
\(Y\). The reference (2) is independent of the homogeneous algebra
generators: it is the unique equivariant function zero on \(K\).

An FS list in A, Definition 2.9 (L277--286), need only generate a
primary ideal. If its homogeneous elements are \(f_j\), then
\(\mathfrak m^r\) is contained in their ideal for some \(r\). At
\(v\in K\), choose \(h_i\) with \(v(h_i)=0\). If all \(v(f_j)>0\), an
expression \(h_i^r=\sum a_j f_j\) has every summand of positive value,
since \(v(a_j)\geq0\), contradicting \(v(h_i^r)=0\). Therefore at least
one \(v(f_j)=0\), while all \(v(f_j)\geq0\). The zero-weight FS function
defined by this primary list is zero on \(K\), and equals (2) everywhere
by (3). This proves reference independence for primary lists as well.

All FS functions factor through \(q\) and satisfy (3). The same is true
of their local uniform limits and decreasing pointwise limits. These are
exactly the CPSH and PSH definitions used here (A, L344--345 and
L762--764), rather than a different intrinsic definition of psh. For a
CPSH function, local uniform convergence on \(Y\) is equivalent to
uniform convergence on \(K\) within a fixed polarization, by (3).

The trivial valuation \(v_{\rm triv}\) is in \(K\). For FS functions, \[
 \sup_K u=u(v_{\rm triv})=\max_j\frac{\lambda_j}{d_\xi(\alpha_j)}.
 \tag{4}
\] This passes to CPSH by uniform approximation. For a decreasing CPSH
net defining \(u\in\operatorname{PSH}(\xi)\), the inequality
\(u(v)\leq u(v_{\rm triv})\) passes pointwise. If the right side were
\(-\infty\), (3) would make \(u\) identically \(-\infty\), which is
excluded. Thus (4), with a finite right side, holds for every actual PSH
function. The normalized representative is \[
 N_\xi u=u-u(v_{\rm triv}).                    \tag{5}
\] It is uniquely characterized by \(\sup_K N_\xi u=0\).

\subsection{2. The orbit-root comparison on the completed
classes}\label{the-orbit-root-comparison-on-the-completed-classes}

For two Reeb fields define \[
 c=c(\xi,\eta)=\min_i\frac{d_\eta(\beta_i)}{d_\xi(\beta_i)},\qquad
 C=C(\xi,\eta)=\max_i\frac{d_\eta(\beta_i)}{d_\xi(\beta_i)}.
 \tag{6}
\] Both are positive. Taking weighted averages in any monomial proves
\(c\leq d_\eta(\alpha)/d_\xi(\alpha)\leq C\) for every supported nonzero
weight \(\alpha\). In particular these constants are independent of
every FS list, its weights, and all approximating potentials.

For \(u\in\operatorname{PSH}(\xi)\) the following inequalities hold for
\(t\geq0\), first termwise for FS functions, then by local uniform and
decreasing limits: \[
 u(v)-Ct\leq u(v+t\eta)\leq u(v)-ct.           \tag{7}
\] The reverse-orbit inequalities follow by applying (7) at \(v-t\eta\).
If \(u(v)\) is finite, (7) makes the function \(t\mapsto u(v+t\eta)\)
finite, strictly decreasing, Lipschitz, and surjective onto
\(\mathbb R\). It therefore has exactly one zero. If \(u(v)=-\infty\),
(7) in both directions makes the entire \(\eta\) orbit \(-\infty\).

Define \(\mathcal R_{\xi\eta}u(v)\) to be this unique zero parameter,
and define it to be \(-\infty\) in the second case. Extend by \(q\) and
put \(\mathcal R_{\xi\eta}u(o)=-\infty\). For an FS presentation this is
the exact formula \[
 u=\max_j\frac{\lambda_j-v(f_j)}{d_\xi(\alpha_j)}
 \quad\Longrightarrow\quad
 \mathcal R_{\xi\eta}u=
 \max_j\frac{\lambda_j-v(f_j)}{d_\eta(\alpha_j)}.       \tag{8}
\] Indeed all terms in \(u(v+t\eta)\) are nonpositive if and only if
\(t\) is at least the right side of (8); at equality one term is zero.
Equation (8) follows from the intrinsic zero definition. It does not
define a map by choosing a presentation. This proves presentation
independence, including redundant generators and primary lists.

\textbf{Theorem 1.} The map \(\mathcal R_{\xi\eta}\) gives order
isomorphisms on \(\mathcal H\), CPSH and PSH. On these classes, \[
 \mathcal R_{\eta\xi}\mathcal R_{\xi\eta}=\mathrm{id},\qquad
 \mathcal R_{\eta\zeta}\mathcal R_{\xi\eta}=\mathcal R_{\xi\zeta}.
 \tag{9}
\] It preserves decreasing nets and their pointwise limits, including
\(-\infty\) behavior. On CPSH it is continuous for local uniform
convergence. On PSH it is a homeomorphism for the topology of pointwise
convergence on \emph{all} analytic points. This last assertion does not
claim continuity for the weaker, quasi-monomial topology used in A,
L805, nor continuity for a strong finite-energy topology.

\emph{Proof.} Formula (8) proves FS membership. If \(u\leq w\), their
zero parameters satisfy \(\mathcal R u\leq\mathcal R w\), because the
orbit functions decrease. If \(u,w\) are CPSH and
\(\|u-w\|_K\leq\epsilon\), equation (3) gives this bound on the entire
source orbit. The lower slope bound in (7) consequently gives \[
 \|\mathcal R u-\mathcal R w\|_K\leq\epsilon/c.       \tag{10}
\] Uniform FS approximation on \(K\), followed by (10), proves that the
intrinsic root lies in CPSH of the target polarization. Its
\(\eta\)-equivariance follows either from the zero definition or from
the approximants; thus convergence on \(K\) implies local uniform
convergence everywhere by the target version of (3).

Now let \(u_a\downarrow u\) be a defining CPSH net. The roots decrease.
If \(r=\mathcal R u(v)\) is finite, \(u(v+(r-\epsilon)\eta)>0\) and
\(u(v+(r+\epsilon)\eta)<0\). Eventually the corresponding values of
\(u_a\) have those signs, so their roots converge to \(r\). If
\(u(v)=-\infty\), then \(u(v-M\eta)=-\infty\) for every \(M>0\).
Eventually \(u_a(v-M\eta)<0\), so their roots are less than \(-M\). This
proves pointwise convergence also in this case. The limit is not
identically \(-\infty\), since \(u(v_{\rm triv})\) is finite. Hence it
belongs to the target PSH by its actual definition.

The root has exactly the same zero set and exactly the same \(-\infty\)
locus as \(u\). Roots along a third Reeb field are therefore identical,
proving the cocycle. Taking the third field to be \(\xi\), where
\(u(v+t\xi)=u(v)-t\), proves the inverse. The inverse is order
preserving, so these are order isomorphisms. The same bracket argument
applies to any pointwise convergent PSH net: at a finite root use two
fixed orbit points, and at \(-\infty\) use \(v-M\eta\). It proves
pointwise continuity of the map and inverse. It also proves preservation
of all existing decreasing PSH limits. Equation (10) proves local
uniform continuity on CPSH. \(\square\)

For completeness, target Reeb variation has local uniform continuity for
a fixed CPSH function. On a compact set of \(v\) the roots remain in a
bounded interval by (7), and \(\eta\) variation in a compact set makes
the associated orbit points range in a compact subset of \(Y\). Uniform
continuity of \(u\) there and the uniform positive lower slope bound
give continuity of the roots. This uses actual continuity of \(u\), not
qualitative rational approximation of energy.

For an FS function, \(u(v_{\rm triv})=0\) means all \(\lambda_j\leq0\)
and at least one \(\lambda_j=0\). The general zero definition gives the
same conclusion about the value at \(v_{\rm triv}\) without needing an
FS presentation: \(\mathcal R u(v_{\rm triv})=0\) if
\(u(v_{\rm triv})=0\). Thus (9) restricts to normalized PSH and induces
genuine bijections of constant quotients by the rule \[
 [u]\longmapsto[\mathcal R_{\xi\eta}(N_\xi u)].       \tag{11}
\] This is independent of representatives because (5) is. The output is
already normalized, so (9) proves inverse and cocycle also for (11).

Normalization before applying the map is essential. For \(R=k[x,y]\),
\(\xi=(1,1)\) and \(\eta=(1,2)\), \[
 \mathcal R_{\xi\eta}(\rho_\xi+2)
 =\max\{-v(x)+2,(-v(y)+2)/2\}.
\] Its difference from \(\rho_\eta\) is \(2\) at \(v_{\rm triv}\) and
\(1\) at the monomial valuation \((v(x),v(y))=(3,0)\in K\). The raw map
does not descend to constant quotients.

For normalized PSH, the exact sign-sensitive comparison on \(K\) is \[
 \frac{u}{c}\leq\mathcal R_{\xi\eta}u\leq\frac{u}{C}\leq0.
 \tag{12}
\] For a nonnegative value of an unnormalized function the inequalities
instead read \(u/C\leq\mathcal R u\leq u/c\). Both follow directly from
(7), including extended values. Let \(Q\) be a compact subset of the
Reeb cone and set \[
 \delta(\xi,\eta)=\max_i\left|\log
 \frac{d_\eta(\beta_i)}{d_\xi(\beta_i)}\right|.
\] Then \(e^{-\delta}\leq c\leq C\leq e^\delta\), and, wherever \(u\) is
finite, \[
 |\mathcal R_{\xi\eta}u-u|\leq(e^\delta-1)(-u)
 \quad\text{on }K.                                  \tag{13}
\] There is \(L_Q<\infty\), depending only on the chosen generators, the
norm on \(N_{\mathbb R}\), and \(Q\), such that
\(\delta\leq L_Q\|\xi-\eta\|\). Indeed the derivative of
\(\log d_\xi(\beta_i)\) has norm at most
\(\|\beta_i\|/\min_{\theta\in Q}d_\theta(\beta_i)\) along the segment;
on the convex hull of \(Q\) this denominator is still bounded below.

For any \emph{fixed} probability measure \(\mu\) on \(K\), (12) also
gives \[
 \int(-\mathcal R u)\mathbf1_{\{-\mathcal R u>A\}}\,d\mu
 \leq c^{-1}\int(-u)\mathbf1_{\{-u>cA\}}\,d\mu.       \tag{14}
\] This is a genuine integrability/tail statement with a fixed measure.
It asserts nothing about moving measures
\(\operatorname{MA}_\eta(\mathcal R u)\).

\subsection{3. Signed finite-type norms and primary-list
representation}\label{signed-finite-type-norms-and-primary-list-representation}

Let \(g_1,\ldots,g_p\) be homogeneous algebra generators, allowing
additional redundant homogeneous elements, with arbitrary real assigned
weights \(\lambda_i\). For \(\alpha\ne0\), define \[
 F^a R_\alpha=\operatorname{span}\{g^{\mathbf e}:
              \text{weight}(g^{\mathbf e})=\alpha,
              \sum e_i\lambda_i\geq a\},\qquad
 \chi(f)=\sup\{a:f\in F^a R_\alpha\}.                \tag{15}
\] For general \(f\) use the minimum of its nonzero homogeneous
components' weights, and give \(1\) weight zero. For fixed \(\alpha\)
there are only finitely many eligible monomials: their exponent sum is
bounded by \(d_\xi(\alpha)/\min_i d_\xi(\operatorname{wt}g_i)\). Thus
(15) is finite, left continuous, exhaustive and separated on each weight
space, and the supremum is attained. It is a quotient non-Archimedean
norm with \[
 \chi(f+g)\geq\min\{\chi(f),\chi(g)\},\qquad
 \chi(fg)\geq\chi(f)+\chi(g).                        \tag{16}
\] These follow by adding and multiplying the spanning expressions. No
assumption that the \(g_i\) are algebraically independent is made. For
homogeneous \(f\), \[
 \min_i\frac{\lambda_i}{d_\xi(\operatorname{wt}g_i)},d_\xi(\alpha)
 \leq\chi(f)\leq
 \max_i\frac{\lambda_i}{d_\xi(\operatorname{wt}g_i)},d_\xi(\alpha).
 \tag{17}
\] Every eligible monomial has weight in that interval, which proves
both inequalities, including the upper bound after quotienting by
relations.

The actual weight \(\chi(g_i)\) may exceed its assigned \(\lambda_i\).
Nevertheless replacing all assigned weights by the actual weights leaves
(15) unchanged. One inequality follows from \(\chi(g_i)\geq\lambda_i\);
the other follows from (16), which bounds the minimum of the actual
monomial weights in every expression by \(\chi(f)\). Taking an
expression attaining the old weight proves attainment also with the
actual weights. This verifies the finite-type condition of B, Definition
3.6 (L493--502), with the saturated norm, rather than assuming that
assigned generator weights are its actual weights.

The exact potential correspondence is \[
 \sup_{0\ne f\in R_\alpha,\ \alpha\ne0}
 \frac{\chi(f)-v(f)}{d_\xi(\alpha)}
 =\max_i\frac{\lambda_i-v(g_i)}{d_\xi(\operatorname{wt}g_i)}.
 \tag{18}
\] To prove the upper bound, take an expression of \(f\) attaining
\(\chi(f)\). The valuation inequality gives
\(\chi(f)-v(f)\leq\max_{\mathbf e}(\sum e_i\lambda_i-v(g^{\mathbf e}))\).
After division by the common degree each monomial expression is a
weighted average of generator expressions, at most the right side of
(18). For the reverse bound use \(f=g_i\) and
\(\chi(g_i)\geq\lambda_i\). The proof permits \(v(g_i)=+\infty\): a
nonzero finite-valued term in an expression supplies the bound, while
entirely infinite-valued expressions give the trivial \(-\infty\) bound.
Thus norm saturation does not alter the FS potential.

An arbitrary primary FS list has a representation of this kind on any
chosen compact Reeb slice \(Q\). Keep its elements \(f_j\) and weights
\(\lambda_j\) and append the fixed algebra generators \(h_i\) with
weight \(-M\). On \(K\) at least one \(f_j\) has value zero, so \[
 u_\xi(v)\geq\min_j\lambda_j/d_\xi(\alpha_j)\geq-B
 \quad(\xi\in Q)
\] for a finite \(B\geq0\). Choose
\(M\geq B\max_{\xi\in Q,i}d_\xi(\beta_i)\). Every appended term is at
most \(-B\) on \(K\), so the potential remains exactly unchanged there.
Equivariance and (3) give exact equality everywhere, for every
\(\xi\in Q\). This is an exact finite-type representation, not an
algebra-generating assumption on the original list. The chosen \(M\) can
depend on this potential; the constants (6), (12), and (22) below do not
depend on \(M\), weights, or potentials.

For rational assigned weights, multiply their filtration indices by a
common denominator. The resulting integer Rees algebra is generated over
\(k[t]\) by the finitely many elements \(g_i t^{-\ell_i}\), where
\(\ell_i\) are the scaled assigned weights. Indeed every filtered
monomial at an index below its assigned weight is its generator product
times a nonnegative power of \(t\); spanning expressions prove the
assertion also in the presence of relations. Localization at \(t\)
recovers \(R[t,t^{-1}]\). Negative weights are allowed. For irrational
weights we assert finite type in the real-norm sense (15), not a
Noetherian integer Rees algebra.

\subsection{4. Valid shifts and uniform signed cutoff
estimates}\label{valid-shifts-and-uniform-signed-cutoff-estimates}

Negating a norm is not a valid treatment of negative weights. For
\(R=k[x]\) with assigned weights \(\lambda_x=\lambda_{x^2}=-1\), (15)
gives \(\chi(x)=-1\) and \(\chi(x^2)=-1>-2\). Then
\(-\chi(x^2)=1<2=(-\chi(x))+(-\chi(x))\), violating multiplicativity.

The valid operation is a degree shift. For any Reeb field \(\beta\) set
\[
 \chi_{A,\beta}(f)=\chi(f)+A d_\beta(\alpha)
 \quad(f\in R_\alpha),                              \tag{19}
\] and extend by minima over homogeneous components. This is precisely
the construction (15) with assigned weights
\(\lambda_i+A d_\beta(\operatorname{wt}g_i)\): every homogeneous
expression acquires the same added weight. It is therefore
multiplicative and finite type. For sufficiently large \(A\), (17)
proves nonnegativity on all of \(R\), including the quotient norm; the
filtration then has \(F^0R=R\) and satisfies every condition of B,
Definition 3.5 (L470--488). It is linearly bounded by (17).

Let \(N_\xi(m)=\dim\bigoplus_{d_\xi(\alpha)\leq m}R_\alpha\),
\(V_\xi=\lim n!N_\xi(m)/m^n>0\), and let \(W_\xi^\chi(m)\) be the sum of
the norm's spectral weights on this direct sum, counting multiplicity.
Spectral weights mean the weights of an adapted basis of the
finite-dimensional filtration, equivalently its jumps with their
dimension multiplicities. They are intrinsic to (15). Define
\(D_\xi^\beta(m)=\sum_{d_\xi(\alpha)\leq m}
d_\beta(\alpha)\dim R_\alpha\). The exact finite identity is \[
 W_\xi^{\chi_{A,\beta}}(m)=W_\xi^\chi(m)+A D_\xi^\beta(m).
 \tag{20}
\] It holds because the entire filtration of \(R_\alpha\) is translated
by one scalar. This is the cutoff identity that must precede subtraction
of limits.

Existence and positivity of \(V_\xi\) are the volume result in B,
Definition 2.4 and Theorem 3.3, L204--209 and L454--460, for the normal
good cone over the algebraically closed characteristic-zero field. Its
displayed \(\mathbb C\) dimensions are ordinary dimensions over the
standing field \(K\) of L181--191. The real-cutoff version follows by
sandwiching between adjacent integer cutoffs. Stieltjes summation gives
the exact identity \(D_\xi^\xi(m)=mN_\xi(m)-\int_0^mN_\xi(t)\,dt\).
Consequently \[
 \lim_{m\to\infty}\frac{D_\xi^\xi(m)}{mN_\xi(m)}=
 \frac{n}{n+1}.                                    \tag{21}
\] To justify the integral limit, divide by \(m^{n+1}\) and use
\(N_\xi(t)=V_\xi t^n/n!+o(t^n)\) beyond an arbitrary fixed cutoff; the
bounded initial integral vanishes after this division.

For each fixed \(\xi\), use (19) with \(\beta=\xi\) to obtain a
nonnegative filtration. B, Corollary 3.13 (L635--645), then gives
convergence of \(W_\xi^{\chi_{A,\xi}}(m)/(mN_\xi(m))\). Its assumptions
have just been checked. Combining the exact identity (20) with the
proved limit (21) proves existence of the signed cutoff mean
\(S_\chi(\xi)=\lim W_\xi^\chi(m)/(mN_\xi(m))\) and the exact rule
\(S_{\chi_{A,\xi}}(\xi)=S_\chi(\xi)+A n/(n+1)\). This is not a
difference of two unproved limits.

If all assigned weights are nonpositive, all spectral weights are
nonpositive by (17). The degree inclusions
\(R_{\xi,\leq m/C}\subset R_{\eta,\leq m}
\subset R_{\xi,\leq m/c}\) give finite inequalities
\(W_\xi^\chi(m/c)\leq W_\eta^\chi(m)\leq W_\xi^\chi(m/C)\). Using the
established limits and volume asymptotics yields \[
 \frac{V_\xi}{V_\eta}c^{-(n+1)}S_\chi(\xi)
 \leq S_\chi(\eta)\leq
 \frac{V_\xi}{V_\eta}C^{-(n+1)}S_\chi(\xi).          \tag{22}
\] The volume inclusions also give
\(V_\xi/C^n\leq V_\eta\leq V_\xi/c^n\). Thus \[
 e^{-(2n+1)\delta}(-S_\chi(\xi))
 \leq -S_\chi(\eta)\leq
 e^{(2n+1)\delta}(-S_\chi(\xi)).                    \tag{23}
\] There is no division by \(S\) in this argument: if it is zero, all
three quantities are zero. These constants are independent of the
assigned weights and appended generators. The fields in (22) may be
irrational and nonproportional.

A normalized primary FS function can be represented by such a
nonpositive norm: replace its weights by
\(\lambda_j-u(v_{\rm triv})d_\xi(\alpha_j)\) and append sufficiently
negative algebra-generator weights. At least one retained weight is
zero. Using the same resulting weights at another Reeb field gives
exactly its normalized orbit comparison. This connects (22) to norm
representations of the correct normalized potentials. It does \emph{not}
identify \(S_\chi\) as a presentation-independent functional of the
potential; nor does it identify \(S_\chi\) with genuine cone energy.

Rational weight approximation is also harmless at this algebraic stage.
If \(|\lambda_i-\lambda_i'|\leq\epsilon d_\xi(\operatorname{wt}g_i)\),
all monomial weights of fixed degree differ by at most
\(\epsilon d_\xi(\alpha)\). Filter inclusions imply that the ordered
spectral weights in \(R_\alpha\) differ by at most that scalar.
Therefore the cutoff sums differ by at most \(\epsilon D_\xi^\xi(m)\),
and the limiting means differ by at most \(\epsilon n/(n+1)\). Equation
(18) also gives uniform potential approximation with error at most
\(\epsilon\). These are norm and potential approximation results; they
do not supply missing analytic energy convergence.

There is a precise cancellation issue. At its own Reeb field, (18) gives
\(u_{\chi_{A,\xi},\xi}=u_{\chi,\xi}+A\), and the genuine energy
translation gives \(E_\xi(u+A)=E_\xi(u)+A\). If a cone representation
were \(E_\xi(u)=a_nS_\chi(\xi)\), its shift coefficient is forced by
(21) to be \(a_n=(n+1)/n\). At another field \(\eta\), a shift by
\(A d_\xi\) changes the numerator by a degree-dependent quantity; it is
not generally a constant shift. One cannot cancel \(A\) at both fields
after using one reference degree. Equations (20)--23) avoid that invalid
operation.

\subsection{5. Genuine normalization checks and an exact projective
result}\label{genuine-normalization-checks-and-an-exact-projective-result}

The genuine affine-cone convention in A, Theorem 2.12 and Definition
4.1, L321--328 and L522--529, is \[
 \operatorname{MA}_\xi(u_1,\ldots,u_{n-1})=
 V_\xi^{-1}dd^c u_1\wedge\cdots\wedge dd^c u_{n-1}
                  \wedge dd^c\rho_\xi^+,
\] \[
 E_\xi(u)=\frac1{nV_\xi}\sum_{j=0}^{n-1}
 \int (u-\rho_\xi)(dd^c u)^j(dd^c\rho_\xi)^{n-1-j}
                              dd^c\rho_\xi^+.       \tag{24}
\] Here \(dd^c=d'd''\) and \(\rho_\xi^+=\max(\rho_\xi,0)\). The mixed
measures have mass one and are carried by \(K\) (A, Theorem 2.12). This
convention has \(n-1\) potential factors and one radial factor, with
\(n\) terms in energy. Substituting \(u=\rho_\xi+a\) in (24), the
currents are unchanged and each integral equals \(aV_\xi\), giving
\(E_\xi(\rho_\xi+a)=a\). This argument uses the actual probability
normalization, not B, Theorem 5.15.

For \(k[x_1,\ldots,x_n]\), \(\xi=(1,\ldots,1)\) and the degree norm
\(\chi(x^{\mathbf e})=|\mathbf e|\), formula (18) gives
\(u=\rho_\xi+1\). For integer \(m\geq1\), \[
 N(m)=\binom{m+n}{n},\qquad
 W(m)=\sum_{|\mathbf e|\leq m}|\mathbf e|
      =n\binom{m+n}{n+1}.
\] The last identity follows by summing
\(d\binom{d+n-1}{n-1}=n\binom{d+n-1}{n}\) and the binomial summation
identity. Thus \(W(m)/(mN(m))=n/(n+1)\) \emph{at every cutoff}, whereas
(24) gives \(E(u)=1\). In particular for \(k[x,y]\) the two values are
\(2/3\) and \(1\). This is a mathematical counterexample to the
unmodified printed equality \(E=S\), not just a numerical diagnostic.

For \(n=1\), on \(k[x]\) with \(d_\xi(x)=a>0\), \(V_\xi=1/a\). On the
calibrated one-dimensional torus skeleton, \(dd^c\max(\log|x|/a,0)\) has
the slope-jump mass \(1/a\) at \(v(x)=0\). Dividing by \(V_\xi\) gives
one Dirac probability mass. The product of zero potential factors in
(24) is the empty product, so \(E(\rho_\xi+c)=c\), while the
degree-shift cutoff coefficient is \(1/2\). This checks the
low-dimensional convention directly. It supplies no nonproportional
comparison in dimension one.

For every cone the same scaling calculation in (24) gives
\(V_{s\xi}=s^{-n}V_\xi\) and \(E_{s\xi}(u/s)=E_\xi(u)/s\), since
\(\rho_{s\xi}=\rho_\xi/s\). There are \(n\) current factors, and the
energy integrand has one more scalar factor. These are proportional
identities only.

There is an exact signed \emph{projective test-configuration}
representation available without a smoothness assumption. Let \(Z\) be
normal projective of dimension \(n-1\), let \(A\) be ample, and let
\((\mathcal Z,\mathcal A)\) be an ample test configuration. Replace
\(A\) by a sufficiently divisible multiple so that its induced
filtration on the full section ring \(\bigoplus_{d\geq0}H^0(Z,dA)\) is
the one in H, Proposition 2.15. Signed rational weights are permitted:
clear denominators by the corresponding integer scaling. By H,
Proposition 5.9 (not a Theorem in this version), the rescaled
degree-wise spectral measures converge to its DH probability measure.
Linear upper growth and the lower bound of H, Remark 5.2, put them in a
common compact interval, so their signed means converge as well. H,
Proposition 3.12(iii), identifies that mean with the actual intersection
number \[
 e=\frac{(\overline{\mathcal A}^{,n})}{n(A^{n-1})}
   =E^{\rm NA}_{\rm proj}(\phi).                    \tag{25}
\] The second equality is H, Definition 7.2. Thus a distribution theorem
is not being substituted for an energy identification.

To check its MA convention, use H, Section 6.7: the mixed model measure
is \(V_A^{-1}\sum_E b_E(\mathcal A_1|_E\cdots
\mathcal A_{n-1}|_E)\delta_{v_E}\), of mass one. If a vertical divisor
has coefficient \(r_E\), its potential value at \(v_E\) is \(r_E/b_E\).
H, Remark 6.12, therefore gives the intersection pairing as \(V_A\)
times its integral against this probability measure. Expanding the
difference of \(n\)-fold intersections using H, Lemma 7.4, proves the
energy formula with \(n\) terms and denominator \(nV_A\). There is no
further \(V_A^{-1}\) after integrating a probability measure. This
derivation fixes the normalization even though the extract of H, Remark
7.5, prints an additional \(1/V\) alongside the statement
\(\mu_j=\operatorname{MA}^{\rm NA}(\cdots)\).

Write \(N_d=h^0(Z,dA)\) and \(W_d\) for the signed total degree-\(d\)
weight. The preceding mean convergence says \(W_d/(dN_d)\to e\). Hilbert
asymptotics give \(N_d=a d^{n-1}+O(d^{n-2})\), with \(a=V_A/(n-1)!\);
for a point (\(n=1\)) use \(N_d=1\) directly. For all \(\epsilon>0\),
beyond a fixed degree \(D\) the error in \(W_d-e dN_d\) is bounded by
\(\epsilon dN_d\). Summing, dividing by \(m\sum_{d\leq m}N_d\), and
making the fixed initial sum negligible proves \[
 S_{\rm cumulative}=\frac{n}{n+1}e,\qquad
 E^{\rm NA}_{\rm proj}(\phi)=\frac{n+1}{n}S_{\rm cumulative}.
 \tag{26}
\] This proof handles signed weights without negating a filtration.
Multiplying the line bundle or clearing weights changes both sides by
the explicitly corresponding scaling; \(V_A\) is always the volume of
the line bundle actually used. Normality of \(Z\) is the hypothesis
here; smoothness is not inferred.

Equations (25)--26) are a proved application to actual projective test
configurations. They are not yet a proof that every primary cone FS
norm, at every rational Reeb field with weighted grading, has exactly
the same energy as (24) after all Veronese and radial calibrations. That
remaining bridge is not hidden in the notation.

\subsection{6. Source audit and the limits of the analytic
argument}\label{source-audit-and-the-limits-of-the-analytic-argument}

\textbf{The energy inconsistency.} B, Lemma 3.16(b), L688--755, states
and proves \(S=\frac{n}{n+1}\widetilde S\), where \(\widetilde S\) is
the degree-wise mean of L679--685. B, Corollary 5.14, L1344--1354,
identifies cone and projective energies under its proposed bridge. The
final proof, L1366--1372, then writes \(E=\widetilde S=S\). The degree
example above disproves this last equality under the printed
definitions. Formula (26) repairs the factor for the projective
situation in its stated scope. It does not prove the universal cone
representation from the diagnostic.

\textbf{Rational quotient and radial calibration.} Normality and
projectivity of a rational quotient can be checked algebraically:
degree-zero localized charts of a normal graded domain are normal, and a
sufficiently divisible Veronese presents \(\operatorname{Proj}R\) as
projective with \(\mathcal O(d)\) ample. But this does not make the
quotient smooth, nor does it prove the radial energy factor. B,
Proposition 5.13, L1282--1304, uses sections of \(mL\), writes a product
trivialization via \(s_0\), and expresses the potential as
\(\pi^*\widetilde\phi-z_0\) in the moment coordinates of those sections.
In those coordinates its preceding expression has radial coefficient
\(-z_0/m\). A nonvanishing section of \(mL\) does not trivialize \(L\).
A valid replacement must account for the degree-\(m\) fiber moment map
and the integral lattice/calibration before cancelling factors. For
\(\eta=s\eta_0\) it must also connect
\(\mathcal O_{\operatorname{Proj}R}(d)/(sd)\), the cone volume, and the
radial current, including ineffective integral actions and nonstandard
weighted gradings. Those operator identifications are not proved here.
F, Theorem 4.13, cannot fill this gap unchanged: F, L525, has the
standing assumption that the projective variety is smooth. H supplies
(25), but still requires the cone-to-model bridge.

\textbf{What CLD does supply.} G, Theorem 15.7.1, L22100--22131, places
a topologically separated pure-dimensional analytic space, with its
analytic boundary, in the calibrated tropical framework. Smoothness is
not a hypothesis; trivial valuation is explicitly included in G,
L585--610. For an analytification of a normal affine variety, separation
and pure dimension hold, and the analytic space without the vertex has
no analytic boundary. Local tropical charts still require their actual
calibration. G, Proposition 8.7.2, L10432--10449 (called a Theorem in
the assignment), permits positive closed currents to vary and gives
weighted-current convergence when all the continuous potentials converge
locally uniformly. The zeroth weight is continuous and locally uniformly
convergent too. It does not handle unbounded decreasing potentials by
itself. G, Corollary 18.3.3(b), L26449--26457, gives the finite-map
factor only when the calibrated generic degree is the specified
constant. Omitting that degree in a Veronese moment-map argument is not
justified.

\textbf{The envelope chain.} E, Definition 3.15(2), L766--768, uses
lower growth \(g\leq u+O(1)\), whereas Definition 6.1(2), L1891--1893,
uses upper growth \(u\leq g+O(1)\). These are different: on
\(\mathbb A^1\), \(u=\log|z|\) is bounded above by \(g=\max(\log|z|,0)\)
and is the decreasing limit of \(\max(\log|z|,-j)\), but is not bounded
below by an FS function up to a constant near \(z=0\). FS functions in E
include a finite constant floor. The coherent convention for this audit
is the upper growth inequality; it is not silently substituted in
quotations.

Changing this inequality does not complete the application. E,
Proposition 4.19, L1560--1576, invokes Lemma 3.25(1) to extend
\(u-cAg\). The invoked lemma, L950--957, explicitly assumes
\(u\in\operatorname{PSH}^{\log}_{\rm reg}\); membership in this
regularizable class is the conclusion of Proposition 4.19. The attempted
chain is Proposition 4.19 to Proposition 6.2(2), then the smooth
starting case and singular induction in Theorem 6.3, then Theorem 6.5,
and finally Theorem 7.1. The complete relevant proof bodies are retained
in the provenance file.

A plausible repair would first establish local truncation/gluing of
\(u-cAg\) on a smooth compactification without regularizability, then
apply a valid smooth local-to-global theorem and global FS
approximation. This would need proof that each truncation is in the
local psh class at every boundary point, the relevant global theorem
under its true net/countability hypotheses, and its propagation through
the singular induction. This document does not supply that repair. It
does not conclude that Theorem 7.1 is false, and does not reject it for
being a v1 preprint. It simply does not use its unverified application
as an indispensable premise.

Theorem 7.1's first stage assumes a countable dense subfield and
descends a psh function from a resolution using connected geometric
fibers and weak-to-global psh comparison. Its final stage, L2718--2747,
descends PL data to countable subfields and uses orthogonality,
base-change of MA and domination to remove the cardinality restriction.
For a trivially valued uncountable field, a countable subfield is not
dense in the discrete field topology. That final stage would therefore
have to be used and justified, not avoided by assuming countability. Any
cone application would further have to check its global reference
metric, the ample rational quotient class after Veronese/scaling,
residual-torus invariance and the identification of the resulting
carrier with this document's fixed \(K\). These are not consequences of
normality alone.

\textbf{Uniform differentiation.} C, Theorem 8.5, L4843--4852, is an
actual projective result for \(E^\downarrow\), with a quadratic
remainder for PL directions, uniform in the potential. The agreement
with enveloped genuine energy for all finite-energy potentials retains
an envelope hypothesis. C, Corollary 8.6, extends differentiability, not
a universal quadratic estimate, to continuous directions. Its proof uses
uniform approximation and the energy's Lipschitz and translation
properties. The two signs are handled by replacing \(f\) by \(-f\). The
extracted proof L4947 and L4954 has mismatched extra factors/signs in
the approximating inequalities: the correct elementary inequality is
\(f_m-\epsilon_m\leq f\leq f_m+\epsilon_m\), which gives the upper
derivative bound with \(+\epsilon_m\), not \(-\epsilon_m\). These
harmless corrections repair that elementary squeeze; they do not
transfer the projective theorem to cones.

An explicit remaining cone target is the following. For a common dense
direction class \(D_Q\subset C(K)\), compact Reeb slice \(Q\), and
normalized energy range \(-E_\xi(u)\leq B\), identify an actual extended
functional \(F_\xi\) with the genuine operator's energy after
enveloping, and prove, for each appropriate dense test direction \(f\),
\[
 |F_\xi(u+tf)-E_\xi(u)-t\int_K f\,\operatorname{MA}_\xi(u)|
 \leq C_{Q,B,f}t^2                                  \tag{27}
\] for both signs and a uniform small-\(t\) interval, or prove another
explicit uniform modulus sufficient for the variational method. No such
\(F_\xi\), envelope agreement, dense class with uniform positivity
domination, or remainder theorem is established here. If (27) were
available only in a dense class, continuous directions would require the
actual uniform-approximation estimate \(2\|f-f_m\|_K+ C_{Q,B,f_m}|t|\)
for the derivative quotient. This implies a uniform differentiability
modulus by choosing \(m\) before \(t\), but gives no quadratic bound for
every continuous \(f\).

D cannot bypass these steps. Its results, including Theorem 4.6, assume
a common linear space of closed forms with mixed products, order
inclusions \(D_\omega\subset D_{\omega'}\), and the associated
integration and positivity framework; its envelope/differentiation
statements also retain orthogonality. The nonlinear root map supplies
none of these mixed-form axioms across different Reeb fields.

\subsection{7. What is still missing from the exact
task}\label{what-is-still-missing-from-the-exact-task}

First, prove an intrinsic representation of (24) for every primary FS
presentation on every rational weighted quotient, including signed
norms, norm saturation, quotient/model normalizations and calibrated
radial factors, and then its genuine CLD continuity to irrational
fields. The projective equality (26) and the translation check fix
necessary factors but do not establish this representation.

Second, turn (22) into an estimate for the \emph{normalized genuine
energy} with constants uniform on \(Q\). Only after this is proved could
a defining decreasing net \(u_a\downarrow u\in\mathcal E^1(\xi)\) be
shown to have target energies bounded below, and hence to define an
element of \(\mathcal E^1(\eta)\). Theorem 1 by itself only preserves
PSH membership.

Third, establish uniform integrability and tails for all required moving
mixed MA measures, for example terms of the form \[
 \int_K(-\mathcal R u_a)\mathbf1_{\{-\mathcal R u_a>A\}}
 \operatorname{MA}_\eta((\mathcal R u_a)^{[j]},
                                     \rho_\eta^{[n-1-j]}),
 \tag{28}
\] and the mixed tests needed when both the potential and polarization
vary. No estimate tending to zero uniformly in the relevant limits is
proved for (28). The fixed-measure bound (14) does not imply one.
Strong-topology continuity and finite-energy mixed-integrability are
consequently unproved.

Fourth, supply the valid envelope/extended-energy construction and the
actual two-sign, Reeb-uniform differentiation estimate (27), with a
proved dense direction class and continuous-direction extension. This is
the missing variational input. Proportional identities, FS
approximation, signed count convergence, and compilation do not
constitute it. The original all-cone, all-Reeb P13.1 therefore remains
unresolved by this draft.

\subsection{8. Evidence and
reproducibility}\label{evidence-and-reproducibility}

The executable \texttt{audit\_sources.py} verifies all eight PDF and
extract hashes (and supplied HTML hashes), preserves LF-indexed literal
statements and proof bodies, and writes
\texttt{source-extracts/source-provenance.json}. It records 51 excerpts
and explicit source-application limits. \texttt{input.json} is unchanged
and \texttt{input-original.json} is byte-identical;
\texttt{assignment-original.txt} preserves the decoded assignment
verbatim, and \texttt{task-original.json} preserves the raw task-object
JSON bytes within the supplied input. The preservation hashes and
environment details are recorded locally.

The light exact-arithmetic program \texttt{checks.py} checks polynomial
degree-count identities, finite orbit-root examples and sign-sensitive
bounds, primary-list supplementation examples, signed shift identities,
and finite cutoff inclusions. Its results in \texttt{checks-output.json}
are finite program observations, not a proof for all cones. The
mathematical proofs are the arguments above. No external compute,
workers, network messages, software installation or other campaign
material is used.

Rerun in this directory:

\begin{Shaded}
\begin{Highlighting}[]
\ExtensionTok{python}\NormalTok{ audit\_sources.py}
\ExtensionTok{python}\NormalTok{ checks.py}
\ExtensionTok{python}\NormalTok{ build.py}
\end{Highlighting}
\end{Shaded}

\texttt{build.py} regenerates the standalone \texttt{P13.1-paper.tex}
from this document using the existing Pandoc executable, then runs the
existing \texttt{pdflatex} twice. Actual stdout and stderr are preserved
separately for conversion and each compilation; each pass's actual TeX
log is copied to a separate filename. The build report records commands,
return codes, versions, hashes and the preservation check. A successful
build establishes readability and reproducibility, not correctness,
novelty, or completion of the mathematical task.

\end{document}
